\documentclass[11pt]{article}
\usepackage{kristiania-workbook}

\begin{document}
\begin{center}
  {\LARGE\bfseries Run a cloud-native workload in Docker}\\[3pt]
  {\large Session 5 -- Cloud-native vs lift-and-shift}\\[5pt]
  {\small Exercise 5 \textperiodcentered\ Session 5 lab \textperiodcentered\ Estimated time: 75--100 min}
\end{center}
\vspace{2pt}\hrule\vspace{1.1em}

Write your answers in the answer document \cmd{SKY2100\_Exercise5\_Answers.docx}. Last week you locked down a host; today you run a real
application as a \textbf{container} and feel the difference between a virtual machine and a container
first-hand. You will keep this Nextcloud workload: in session~6 you put it behind HTTPS, and in
session~7 you configure its access control.

\noindent\textbf{Caution:} Run everything in your own VM. Containers here are for learning; do not
expose this Nextcloud to the public internet.

\wbsection{1\quad Install Docker}
\task{Install and verify Docker} (note the version in the answer document):
\begin{quote}\ttfamily
sudo apt update \&\& sudo apt install -y docker.io\\
sudo docker --version\\
sudo docker run hello-world
\end{quote}

\wbsection{2\quad Run Nextcloud as a container}
\task{Pull and start the container:}
\begin{quote}\ttfamily
sudo docker run -d -p 8080:80 --name nextcloud nextcloud
\end{quote}
Open \cmd{http://localhost:8080} in the VM browser and complete the first-run admin setup.

\task{Inspect the running container:} \cmd{sudo docker ps}, \cmd{sudo docker logs nextcloud}.

\wbsection{3\quad Container vs VM --- observe the difference}
\task{Note in the answer document, and compare with your VM baseline (Exercise 2):} the approximate
start-up time of the container, the ports it exposes (\cmd{docker ps}), and whether it runs its own
kernel (Y/N).
\task{In one or two sentences: why did the container start so much faster than booting a VM?}

\aha{The container started in seconds because it did not boot an operating system. Your VM from
Exercise~2 boots a whole kernel and its background services; the Nextcloud container reuses the host's
kernel and is really just an isolated process. That is why it is small and fast, and it is also why its
isolation is weaker than a VM's. A VM is separated by the hypervisor, but every container on this host
shares one kernel, so a kernel-level escape reaches the host and its neighbours. The speed you liked
and the security you have to think about come from the same design choice.}

\wbsection{4\quad Cloud transfer: Microsoft Learn}
\task{Complete ``Introduction to Docker containers''.}
\par\noindent\mbox{\href{https://learn.microsoft.com/training/modules/intro-to-docker-containers/}{learn.microsoft.com/training/modules/intro-to-docker-containers}}

\task{Define the terms} you met in the lab, one line each in the answer document: image, container,
registry (Docker Hub).

\wbsection{5\quad Azure reflection}
\task{You ran Nextcloud with \cmd{docker run}. Azure Container Apps would run the \emph{same} image
for you. Name one thing you would no longer manage, and one new thing you would have to configure.}

\wbsection{6\quad Knowledge check}
\begin{enumerate}
\item Order these from least to most cloud change: refactor, rehost, rebuild, replatform.
\hint{Rehost is the least change; rebuild is the most.}
\item What is the key technical difference between a \textbf{VM} and a \textbf{container}?
\hint{One runs its own kernel; one shares the host's.}
\item Which kernel does a container use?\\
      \cbox its own \quad \cbox a copy of the host's \quad \cbox the shared host kernel \quad \cbox none
\hint{It shares the host's.}
\item Name the two operating-system features that make containers possible (Comer Ch.~6).
\hint{Namespaces give isolation, control groups (cgroups) impose limits.}
\item Define \textbf{image}, \textbf{container} and \textbf{Dockerfile} in one line each.
\hint{A read-only template, a running instance of it, and the recipe that builds the template.}
\item Give \emph{two} container-specific security risks that do not apply to a single VM.
\hint{Shared-kernel escape, untrusted or outdated images, secrets baked into a layer.}
\item Contrast a \textbf{monolith} with \textbf{microservices} in one sentence each.
\hint{One codebase with everything, versus many small independent services.}
\item State one \emph{advantage} and one \emph{disadvantage} of microservices (Comer Ch.~12).
\hint{Independent scaling and fault containment — at the cost of many network calls.}
\item What does a \textbf{service mesh} do, and which security control can it enforce on internal
      calls?
\hint{A sidecar proxy beside each service; it can enforce mutual TLS.}
\item As you move from lift-and-shift to serverless, who becomes responsible for patching the
      \textbf{operating system}?
\hint{As you move toward serverless, the provider takes it over.}
\item True or false: cloud-native architecture removes security risk. Justify.
\hint{It relocates risk (to configuration, images, identity); it does not remove it.}
\item (Reflection) For your Nextcloud, name one workload-security control (CSA Domain~8) you would
      apply first, and why.
\end{enumerate}

\signoff

\clearpage
\solheader
{\LARGE\bfseries\color{kred} Solutions}\\[2pt]
{\small Work through the lab first, then check here. Model answers are indicative — equivalent reasoning is fine.}\par\vspace{8pt}
\noindent\textbf{Note:} Model answers are indicative -- accept equivalent reasoning. Multiple-choice
answers are in \textbf{bold}.

\wbsection{Knowledge check}
\begin{enumerate}
\item \textbf{Rehost $\rightarrow$ replatform $\rightarrow$ refactor $\rightarrow$ rebuild.}
\item A \textbf{VM} virtualises the \emph{hardware} and runs its own OS kernel; a \textbf{container}
      virtualises the \emph{operating system} and shares the host kernel -- it is essentially an
      isolated process.
\item \textbf{The shared host kernel.}
\item \textbf{Namespaces} (give a process its own view of processes, network, filesystem, users) and
      \textbf{control groups / cgroups} (limit CPU, memory, I/O). (Comer Ch.~6, p.~73--74, verified.)
\item \textbf{Image} = a read-only template (app + dependencies, built in layers); \textbf{Container} =
      a running instance of an image; \textbf{Dockerfile} = a recipe that builds an image step by step.
\item Any two of: shared-kernel escape (weaker isolation than a VM); untrusted/outdated images;
      secrets baked into an image layer; over-privileged or root containers / host networking.
\item \textbf{Monolith} = one large program with all features in a single codebase; \textbf{microservices}
      = many small independent services communicating over the network. (Comer Ch.~12, p.~163--165.)
\item Advantage: independent scaling/deployment, fault containment, parallel teams. Disadvantage: many
      network calls (complexity/latency), harder to test and trace, more connections to secure. (Comer
      Ch.~12, p.~165--167, verified.)
\item A \textbf{service mesh} puts a sidecar proxy beside each service to manage service-to-service
      traffic; it can enforce \textbf{mutual TLS (mTLS)}, identity and policy on internal calls. (Comer
      Ch.~12, p.~175.)
\item \textbf{The provider} (as you move toward serverless / PaaS).
\item \textbf{False} -- cloud-native \emph{relocates} risk (from patching servers to securing
      configuration, images and identity); it does not remove it.
\item Reflection -- e.g.\ harden the image, run non-root, scan for vulnerabilities, or segment the
      workload (CSA Domain~8).
\end{enumerate}

\wbsection{Lab checkpoints}
\begin{itemize}
\item The container starts in seconds because it does \emph{not} boot an operating system -- it is a
      process using the host kernel (contrast with the VM boot in Exercise~2).
\item \cmd{docker ps} should show the \cmd{nextcloud} container mapping host \cmd{8080} to container
      \cmd{80}. The container does \emph{not} run its own kernel.
\item Good answers connect ``faster start-up'' to ``no guest OS / shared kernel'', not just ``smaller''.
\end{itemize}

\wbsection{References}
Comer Ch.~6 -- Containers (VM overhead \textbf{p.~71--72}, namespaces \textbf{p.~74}, container approach
\textbf{p.~75}, Docker images/Dockerfile \textbf{p.~76--83}); Comer Ch.~12 -- Microservices (monolith vs
microservices \textbf{p.~163--165}, advantages/disadvantages \textbf{p.~165--167}, service mesh
\textbf{p.~175}) -- all verified. CSA Security Guidance v5, \textbf{Domain~8: Cloud Workload Security}
(verified). J\o sang (application/workload security -- syllabus, page numbers not verified). MS Learn:
\emph{Introduction to Docker containers} (verified, June 2026).

\end{document}
